% TODO make hw stylesheet
\documentclass[10pt]{article}
\usepackage[letterpaper,top=1in]{geometry}
\usepackage[dvipsnames,svgnames]{xcolor}
\usepackage[shortlabels]{enumitem}
\usepackage[framemethod=TikZ]{mdframed}
\usepackage{amsmath,amssymb,amsthm}
\usepackage[colorlinks]{hyperref}
\usepackage{multicol}
\usepackage{tabularx}

\hypersetup{
	colorlinks=true,
	linkcolor=blue,
	filecolor=magenta,
	urlcolor=magenta,
	pdftitle={MAT 528 HW}
}

\usepackage{mathtools}
\usepackage{thmtools}
\usepackage[textsize=footnotesize]{todonotes}
\usepackage{titling}

\reversemarginpar
\mdfdefinestyle{mdbluebox}{roundcorner=10pt,innerbottommargin=9pt,
	linecolor=blue,backgroundcolor=TealBlue!5,}
\declaretheoremstyle[headfont=\sffamily\bfseries\color{MidnightBlue},
mdframed={style=mdbluebox},]{thmbluebox}

\mdfdefinestyle{mdbloobox}{frametitlefont=\bfseries,innerbottommargin=8pt,
	nobreak=true,backgroundcolor=TealBlue!5,linecolor=blue,}
\declaretheoremstyle[headfont=\bfseries\color{MidnightBlue},
mdframed={style=mdbloobox},headpunct={\\[3pt]},postheadspace=0pt,]{thmbloobox}

\mdfdefinestyle{mdredbox}{frametitlefont=\bfseries,innerbottommargin=8pt,
	nobreak=true,backgroundcolor=Salmon!5,linecolor=RawSienna,}
\declaretheoremstyle[headfont=\bfseries\color{RawSienna},
mdframed={style=mdredbox},headpunct={\\[3pt]},postheadspace=0pt,]{thmredbox}

\mdfdefinestyle{mdgreenbox}{linecolor=ForestGreen,backgroundcolor=ForestGreen!5,
	linewidth=2pt,rightline=false,leftline=true,topline=false,bottomline=false,}
\declaretheoremstyle[headfont=\bfseries\sffamily\color{ForestGreen!70!black},
mdframed={style=mdgreenbox},headpunct={ --- },]{thmgreenbox}

\mdfdefinestyle{mdorangebox}{linecolor=RedOrange,backgroundcolor=RedOrange!5,
	linewidth=2pt,rightline=false,leftline=true,topline=false,bottomline=false,}
\declaretheoremstyle[headfont=\bfseries\sffamily\color{RedOrange!70!black},
mdframed={style=mdorangebox},headpunct={ --- },]{thmorangebox}

\mdfdefinestyle{mdblackbox}{linecolor=black,backgroundcolor=RedViolet!5!gray!5,
	linewidth=3pt,rightline=false,leftline=true,topline=false,bottomline=false,}
\declaretheoremstyle[mdframed={style=mdblackbox}]{thmblackbox}

\declaretheoremstyle[spaceabove=6pt,spacebelow=6pt]{basehead}

\declaretheorem[style=basehead,name=Problem]{problem}
\declaretheorem[style=thmredbox,name=Problem,numbered=no]{problem*}
\declaretheorem[style=thmbloobox,name=Setup]{setup} % for config geo
\declaretheorem[style=thmredbox,name=Example,sibling=problem]{example}
\declaretheorem[style=thmredbox,name=$\bigstar$ Problem,sibling=problem]{niceprob}
\declaretheorem[style=thmbluebox,name=Theorem,sibling=problem]{theorem}
\declaretheorem[style=thmbluebox,name=Corollary,sibling=theorem]{corollary}
\declaretheorem[style=thmbluebox,name=Proposition,sibling=theorem]{prop}
\declaretheorem[style=thmbluebox,name=Lemma,numberwithin=problem]{lemma}
\declaretheorem[style=thmbluebox,name=Theorem,numbered=no]{theorem*}
\declaretheorem[style=thmbluebox,name=Lemma,numbered=no]{lemma*}
\declaretheorem[style=thmgreenbox,name=Claim,numberwithin=problem]{claim}
\declaretheorem[style=thmgreenbox,name=Claim,numbered=no]{claim*}
\declaretheorem[style=thmblackbox,name=Remark,numberwithin=problem]{remark}
\declaretheorem[style=thmblackbox,name=Remark,numbered=no]{remark*}
\declaretheorem[style=thmgreenbox,name=Definition,sibling=theorem]{definition}
\declaretheorem[style=thmgreenbox,name=Definition,numbered=no]{definition*}

\providecommand{\ol}{\overline}
\providecommand{\ceil}[1]{\left\lceil #1 \right\rceil}
\providecommand{\floor}[1]{\left\lfloor #1 \right\rfloor}
\newcommand{\dd}{\mathrm{d}}
\providecommand{\ul}{\underline}
\providecommand{\eps}{\varepsilon}
\providecommand{\half}{\frac{1}{2}}
\providecommand{\CC}{\mathbb C}
\providecommand{\FF}{\mathbb F}
\providecommand{\NN}{\mathbb N}
\providecommand{\QQ}{\mathbb Q}
\providecommand{\RR}{\mathbb R}
\providecommand{\ZZ}{\mathbb Z}
\providecommand{\dg}{^\circ}
\providecommand{\ii}{\item}
\providecommand{\setdiff}{\smallsetminus}
\providecommand{\alert}[1]{{\sffamily\textbf{\textcolor{blue}{#1}}}}
\providecommand{\inv}{^{-1}}
\providecommand{\mb}[1]{\mathbf{#1}}

\newenvironment{soln}{\begin{proof}[Solution]}{\end{proof}}

\providecommand{\pow}{\operatorname{Pow}}
\providecommand{\vocab}[1]{{\textbf{\textcolor{ForestGreen}{#1}}}}
\providecommand{\lcm}{\operatorname{lcm}}
\providecommand{\abs}[1]{\left|#1\right|}

\usepackage{mathrsfs}

% place title at top right
\setlength{\droptitle}{-8ex}
\pretitle{\begin{flushright}\Large\bfseries}
\posttitle{\par\end{flushright}}
\preauthor{\begin{flushright}}
\postauthor{\end{flushright}}
\predate{\begin{flushright}}
\postdate{\end{flushright}}

\title{MAT 528 HW06}
\author{\large Aditya Pahuja}
\date{\large Due 10/09}
\begin{document}
\maketitle

\begin{problem}
    [Munkres 26.3]
    Show that a finite union of compact subspaces of $X$ is compact.
\end{problem}

\begin{soln}
    Let $C_1$, \dots, $C_n$ be compact subspaces of $X$.
    If $\mathcal U = \{U_\alpha\}_{\alpha\in I}$ is an open cover
    of $C_1\cup\cdots\cup C_n$, then it is an open cover
    of each $C_k$ for $k = 1,\dots,n$. Since $C_k$ is compact,
    it has a finite subcover $\mathcal U_k\subseteq \mathcal U$,
    so $C_1\cup\cdots\cup C_n$ is has the finite open cover
    $\mathcal U_1\cup\cdots\cup\mathcal U_n\subseteq\mathcal U$.
    Thus, given any open cover of $C_1\cup\cdots\cup C_n$
    we can extract a finite subcover,
    which is what it means for $C_1\cup\cdots\cup C_n$ to be compact.
\end{soln}

\begin{problem}
    [Munkres 26.4]
    Show that every compact subspace of a metric space
    is bounded in that metric space. Find a metric space
    in which not every closed bounded space is compact.
\end{problem}

\begin{soln}
    Suppose $C$ is a compact subspace of metric space $X$ with metric $d$.
    Then, fixing $\eps > 0$, $C$ has an open cover
    $\{B_\eps(x)\}_{x\in C}$ consisting of balls centered at each $x\in C$,
    and this open cover has a finite subcover
    $\mathcal U = \{B_\eps(x_j) : x_1,\dots,x_n\in C\}$
    of open balls corresponding to some finitely many points $x_1,\dots,x_n\in C$.

    If $x_0\in X$ is an arbitrary fixed point, then for any $y\in X$,
    $B_{d(x_0,y) + \eps}(x_0)$ contains $B_\eps(y)$ because
    if $y'\in B_\eps(y)$ then
    \begin{equation*}
        d(x_0,y')\le d(x_0,y) + d(y,y') < \eps.
    \end{equation*}
    This means that for $k=1,\dots,n$,
    $B_\eps(x_k)\in\mathcal U$ is contained in $B_{d(x_0,x_k)+\eps}(x_0)$, so
    if $M = \max\{d(x_0,x_k) : k=1,\dots,n\}$ then
    \begin{equation*}
	C\subseteq \bigcup_{k=1}^n B_\eps(x_k) = B_{M+\eps}(x_0),
    \end{equation*}
    so $C$ is bounded because it is contained in an open ball.
    (more specifically, $C$ is bounded because
    $d(x,y) < 2(M+\eps)$ for any $x,y\in B_{M+\eps}(x_0)$,
    hence the same inequality holds for any $x,y\in C$,
    so $C$ has diameter at most $2(M+\eps)$)
    \\

    On the other hand, give $\RR$ the metric
    \begin{align*}
	d(x,y) = 1\text{ if }x\ne y\\
	d(x,y) = 0\text{ if }x=y.
    \end{align*}
    Then $\RR$ is closed (by definition) and bounded (because $d(x,y) \le 1$
    for all $x,y\in\RR$) but $\RR$ is not compact
    because, by example 1 of section 20 in Munkres,
    $d$ induces the discrete topology, so the collection of singleton sets
    $\{\{x\}\}_{x\in\RR}$ is an open cover, which cannot have a finite subcover
    since the union of any finite subcollection of these singleton sets
    is a finite set, and $\RR$ is not finite.
\end{soln}

\begin{problem}
    [Munkres 26.7]
    Show that if $Y$ is compact, then the projection
    $\pi_1\colon X\times Y\to X$ is a closed map.
\end{problem}

\begin{soln}
    Let $F\subseteq X\times Y$ be a closed set.
    We show that the complement of $\pi_1(F)$ is open.
    Let $x_0\in X\setminus\pi_1(F)$;
    to show that $X\setminus\pi_1(F)$ we need to show that
    there exists an open set $W$ such that $x_0\in W\subseteq X\setminus\pi_1(F)$.
    The set $(X\times Y)\setminus F$
    is an open set containing $\{x_0\}\times Y$,
    so by the tube lemma $(X\times Y)\setminus F$ contains
    a tube $W\times Y$ where $W\subseteq X$ is an open set containing $x_0$.
    In particular, no point of $W\times Y$ is in $F$, so
    $\pi_1(W)\subseteq X\setminus\pi_1(F)$ is an open set in $X$ (because
    $\pi_1$ is an open map, as shown on a previous homework)
    that contains $x_0$; this proves that $X\setminus\pi_1(F)$ is open.
\end{soln}

\begin{problem}
    [Munkres 26.8]
    Let $Y$ be a compact Hausdorff space
    and let $f\colon X\to Y$. Then $f$ is continuous
    if and only if the \vocab{graph} $G_f$ of $f$,
    \begin{equation*}
	G_f \coloneq \{(x,f(x)) : x\in X\},
    \end{equation*}
    is closed in $X\times Y$.
\end{problem}

\begin{soln}
    Suppose $f$ is continuous and let $(x_0,y_0)$ be
    a point in $(X\times Y)\setminus G_f$.
    Then $f(x_0)\ne y_0$, so there exist disjoint open sets
    $U,V\subseteq Y$ such that $f(x_0)\in U$ and $y_0\in V$
    because $Y$ is Hausdorff. Then, the open set
    $f\inv(U)\times V$ is disjoint from $G_f$
    because $V$ does not contain any points in $f(f\inv(U)) = U$,
    which is enough to prove that the complement of $G_f$ is open in $X\times Y$,
    hence $G_f$ is closed.

    Conversely suppose that $G_f$ is closed in $X\times Y$.
    Let $V$ be an open set in $Y$ and let $F = (X\times V)\cap G_f$,
    which is closed. Then $F$ is the set
    \begin{equation*}
	\{(x,y)\in G_f : y\in V\} =
	\{(x,f(x))\in X\times Y : x\in X\text{ such that }f(x)\in V\},
    \end{equation*}
    so $\pi_1(F) = f\inv(V)$ because it is the set of all $x\in X$
    for which $f(x)\in V$; since $Y$ is compact, the preceding problem
    implies that $\pi_1$ is a closed map, so $\pi_1(F) = f\inv(V)$ is closed
    for every closed set $F$, implying $F$ continuous.
\end{soln}

\begin{problem}
    [Munkres 27.2ab]
    Let $X$ be a metric space with metric $d$ and let $A\subseteq X$ be nonempty.
    \begin{enumerate}[(a)]
        \ii Show that $d(x,A) = 0$ if and only if $x\in\ol A$.
	\ii Show that if $A$ is compact, $d(x,A) = d(x,a)$ for some $a\in A$.
    \end{enumerate}
\end{problem}

\begin{soln}
    (a) If $x\in\ol A$ then for each $\eps>0$
    the ball $B_\eps(x)$ contains some point $y\in A$, so
    $d(x,A) \le d(x,y) < \eps$; this implies $d(x,A) = 0$.

    Conversely, if $X\notin\ol A$ then there exists $\eps>0$
    such that $B_\eps(x)$ does not intersect $A$, meaning
    for any $y\in A$, $d(x,y)\ge \eps$; taking the infimum over all $y\in A$
    gives $d(x,A) \ge \eps > 0$.\\

    (b) The function $d(x,-)\colon X\to\RR$ is continuous at each $y\in X$ because
    for each $\eps>0$,
    \begin{equation*}
	d(y,y') < \eps \implies \abs{d(x,y) - d(x,y')} \le d(y,y') < \eps
    \end{equation*}
    where $\abs{d(x,y) - d(x,y')} \le d(y,y')$ because
    $d(x,y) \le d(x,y') + d(y,y')$ and $d(x,y') \le d(x,y) + d(y',y)$.
    Therefore by Munkres's Theorem 27.4, since $A$ is compact,
    $d(x,-)$ there exists a point $a\in A$ for which
    \begin{equation*}
	d(x,a) = \inf_{y\in A}d(x,y) = d(x,A).\qedhere
    \end{equation*}
\end{soln}

\begin{problem}
    [Munkres 27.6]
    Let $A_0$ be the closed interval $[0,1]$ in $\RR$.
    Let $A_1$ be the set obtained from $A_0$ by deleting its ``middle third''
    $(\frac13,\frac23)$. Let $A_2$ be the set obtained from $A_1$
    by deleting its ``middle thirds'' $(\frac19,\frac29)$ and $(\frac79,\frac89)$.
    In general, define $A_n$ by the equation
    \begin{equation*}
	A_n = A_{n-1}\setminus\bigcup_{k=0}^\infty
	\left(\frac{1+3k}{3^n},\frac{2+3k}{3^n}\right).
    \end{equation*}
    The intersection $C = \bigcap_{n=1}^\infty A_n$ is called the
    \vocab{Cantor set}; it is a subspace of $[0,1]$.
    \begin{enumerate}[(a)]
        \ii Show that $C$ is totally disconnected.
	\ii Show that $C$ is compact.
	\ii Show that each set $A_n$ is a union of
	finitely many disjoint closed intervals of length $3^{-n}$,
	and show that the endpoints of these intervals lie in $C$.
	\ii Show that $C$ has no isolated points.
	\ii Conclude that $C$ is uncountable.
    \end{enumerate}
\end{problem}

\begin{soln}
    (a) Suppose $A\subseteq C$ contains two distinct points $x,y$.
    Then, pick $n$ so that $3^{-n} < \abs{x-y}$, so that by (c),
    $x$ and $y$ must lie in distinct intervals of $A_n$.
    If $I$ is the interval containing $x$ then $I$ and $A_n\setminus I$
    are both closed in $A_n$ because they are unions of finitely many closed intervals,
    nonempty, and disjoint, so $I\cap A$ and $A\cap(A_n\setminus I) = A\setminus I$
    are closed in $I$, nonempty (since $I$ contains $x$ but not $y$
    so $A\setminus I$ contains $y$), and disjoint, which means
    we have found a separation of $A$.
    Thus $C$ is totally disconnected
    because we have shown that any set $A\subseteq C$
    with more than two elements is disconnected.
    \\

    (b) Each $A_n$ is closed because it is the intersection of closed sets
    \begin{equation*}
	A_n = A_{n-1}\setminus\bigcup_{k=0}^\infty
	\left(\frac{1+3k}{3^n},\frac{2+3k}{3^n}\right)
	= A_{n-1}\cap\underbrace{\left(\RR\setminus\bigcup_{k=0}^\infty
	\left(\frac{1+3k}{3^n},\frac{2+3k}{3^n}\right)\right)}_
	{\text{complement of open, so closed}}.
    \end{equation*}
    Thus the intersection $C = \cap_{n=1}^\infty A_n$ of closed sets
    is closed, and $C$ is bounded because it is contained in the bounded set $[0,1]$,
    which means it is compact by Munkres's Theorem 27.3.
    \\

    (c) This follows from induction: clearly it is true for $A_0$,
    and if $A_n$ consists of finitely many closed intervals of length $3^{-n}$,
    then removing the middle third $((2a+b)/3, (a+2b)/3)$ of
    some closed interval $[a,b]$ leaves behind the closed intervals
    $[a,(2a+b)/3]$, $[(a+2b)/3, b]$ with a third of the length.
    Furthermore, since removing the middle third
    necessarily means that only interior points of each interval are removed
    and the endpoints of each interval stay as endpoints of some new intervals
    (as in the one-interval example $[a,b]$ above), one can never lose the endpoints
    via the middle-thirds-deletion process, so the endpoints are in $C$.
    \\

    (d) For each $x\in C$ and $n\in\NN$, since $x\in A_n$
    it is contained in some interval of length $3^{-n}$ that is contained in $A_n$.
    The endpoints of this interval are in $C$ by part (c)
    and are at most $3^{-n}$ away from $x$; let $x_n$ be one of these endpoints
    (if $x$ is an endpoint of the interval then pick $x_n$ to be the other endpoint).
    We have thus found a point $x_n\in C$ not equal to $x$
    such that $\abs{x-x_n} < 3^{-n}$. This means that $\lim_{n\to\infty}x_n = x$,
    so any open set around $x$ contains some (in fact cofinitely many)
    $x_n$ (none of which are equal to $x$), so $x$ is not an isolated point.
    \\

    (e) Since $C$ is a nonempty compact Hausdorff space with no isolated points,
    Munkres's Theorem 27.7 implies that $C$ is uncountable.
\end{soln}










\end{document}

